\documentclass[10pt,a4paper]{amsart}
\usepackage{amsmath,amssymb,amsthm,graphicx}
\usepackage{verbatim}
\newenvironment{reference}{\comment}{\endcomment}
%\newenvironment{reference}{}{}
\newenvironment{slogan}{\comment}{\endcomment}
\newenvironment{history}{\comment}{\endcomment}

% For commutative diagrams we use Xy-pic
\usepackage[all]{xy}

% We use 2cell for 2-commutative diagrams.
\xyoption{2cell}
% Theorem environments.
%
\theoremstyle{plain}
\newtheorem{theorem}[subsection]{Theorem}
\newtheorem{proposition}[subsection]{Proposition}
\newtheorem{lemma}[subsection]{Lemma}

\theoremstyle{definition}
\newtheorem{definition}[subsection]{Definition}
\newtheorem{example}[subsection]{Example}
\newtheorem{exercise}[subsection]{Exercise}
\newtheorem{situation}[subsection]{Situation}

\theoremstyle{remark}
\newtheorem{remark}[subsection]{Remark}
\newtheorem{remarks}[subsection]{Remarks}

\numberwithin{equation}{subsection}

% Macros
%
\def\lim{\mathop{\mathrm{lim}}\nolimits}
\def\colim{\mathop{\mathrm{colim}}\nolimits}
\def\Spec{\mathop{\mathrm{Spec}}}
\def\Hom{\mathop{\mathrm{Hom}}\nolimits}
\def\Ext{\mathop{\mathrm{Ext}}\nolimits}
\def\SheafHom{\mathop{\mathcal{H}\!\mathit{om}}\nolimits}
\def\SheafExt{\mathop{\mathcal{E}\!\mathit{xt}}\nolimits}
\def\Sch{\mathit{Sch}}
\def\Mor{\mathop{\mathrm{Mor}}\nolimits}
\def\Ob{\mathop{\mathrm{Ob}}\nolimits}
\def\Sh{\mathop{\mathit{Sh}}\nolimits}
\def\NL{\mathop{N\!L}\nolimits}
\def\CH{\mathop{\mathrm{CH}}\nolimits}
\def\proetale{{pro\text{-}\acute{e}tale}}
\def\etale{{\acute{e}tale}}
\def\QCoh{\mathit{QCoh}}
\def\Ker{\mathop{\mathrm{Ker}}}
\def\Im{\mathop{\mathrm{Im}}}
\def\Coker{\mathop{\mathrm{Coker}}}
\def\Coim{\mathop{\mathrm{Coim}}}

% Boxtimes
%
\DeclareMathSymbol{\boxtimes}{\mathbin}{AMSa}{"02}

%
% Macros for moduli stacks/spaces
%
\def\QCohstack{\mathcal{QC}\!\mathit{oh}}
\def\Cohstack{\mathcal{C}\!\mathit{oh}}
\def\Spacesstack{\mathcal{S}\!\mathit{paces}}
\def\Quotfunctor{\mathrm{Quot}}
\def\Hilbfunctor{\mathrm{Hilb}}
\def\Curvesstack{\mathcal{C}\!\mathit{urves}}
\def\Polarizedstack{\mathcal{P}\!\mathit{olarized}}
\def\Complexesstack{\mathcal{C}\!\mathit{omplexes}}
% \Pic is the operator that assigns to X its picard group, usage \Pic(X)
% \Picardstack_{X/B} denotes the Picard stack of X over B
% \Picardfunctor_{X/B} denotes the Picard functor of X over B
\def\Pic{\mathop{\mathrm{Pic}}\nolimits}
\def\Picardstack{\mathcal{P}\!\mathit{ic}}
\def\Picardfunctor{\mathrm{Pic}}
\def\Deformationcategory{\mathcal{D}\!\mathit{ef}}
\usepackage{cleveref}
\crefname{originallemma}{original lemma}{original lemmas}
\Crefname{originallemma}{Original lemma}{Original lemmas}
\makeatletter
\newlabel{lemma-NAK}{{00DV}{source}{}{sourcepointer-1}{}}
\newlabel{lemma-NAK@cref}{{[originallemma][0][]00DV}{[1][0][]source}{}{}{}}
\makeatother
\title{A submodule generating a finite free module over a semi-local ring contains a basis}
\author{The Stacks Project Authors}
\date{stacks-project:algebra.tex:local-snapshot-20261012; GNU FDL 1.2 or later}
\begin{document}
\maketitle
\noindent\textit{Editorial note.} All delivered mathematical text is the author's own, copied from the pinned Stacks Project chapter file algebra.tex; no statement or proof step was rewritten, reordered or removed.\par\medskip
\noindent\textit{Editorial note.} Preamble provenance, as required for this source: the frozen snapshot stages the chapter's own preamble file beside it (preamble.tex, 217 lines, 8116 bytes, the same shared Stacks-wide preamble the published chow and etale-cohomology cards use); the first line of the pinned chapter file is \textbackslash{}input\{preamble\}, so no preamble text lives inside the chapter itself. The card takes that file's lines 8-18 (the verbatim package and the three \textbackslash{}comment-based environments reference/slogan/history at lines 9-12, section 13.2, together with the Xy-pic setup) and its lines 154 to the runtime end of the file (the native theorem declarations and mathematical macros). Not one character was changed, no package was added, no macro was defined, and the excerpt serves only as preamble, never as a span coordinate.\par\medskip
\noindent\textit{Editorial note.} The same-chapter pointers are located at their \textbackslash{}label lines inside the pinned chapter file; the cross-chapter pointers are located in the project's own tag list staged in the snapshot (sources/stacks-chow/tags.txt), at the verbatim lines carrying their canonical tags.\par\medskip
\noindent\textit{Editorial note.} The counted proof refers to Nakayama's lemma by label; it is recorded as a navigation stub pointed at its \textbackslash{}label line, and the argument consuming it stays inside the counted span.\par\medskip
\noindent\textit{Editorial note.} The counted spans contain no \textbackslash{}cite, so the card carries no bibliography.\par\medskip
\noindent\textit{Allowed standard prerequisites (editorial).}\begin{itemize}
\item Nakayama's lemma. Invoked in the counted proof at line 19545. Reduces the search for a basis to a generating condition modulo the Jacobson radical.
\item Linear combinations over an infinite field. Used in the counted proof at lines 19556-19558. Produces a single element of the submodule with a nonzero component in every local factor.
\end{itemize}\medskip
\ifdefined\hypertarget\hypertarget{sourcepointer-1}{}\fi
\noindent\textit{Source reference.} Original lemma 00DV: sources/stacks-algebra/pinned-original-source.tex, source lines 3833--3833. Navigation pointer for \textbackslash{}ref\{lemma-NAK\} used inside the counted span (line 19545); the argument that consumes it is retained in the span.\par\medskip
\begin{lemma}
\label{lemma-semi-local-module-basis-in-submodule}
Let $R$ be a local ring with maximal ideal $\mathfrak m$ and
infinite residue field.
Let $R \to S$ be a ring map.
Let $M$ be an $S$-module and let $N \subset M$ be an $R$-submodule.
Assume
\begin{enumerate}
\item $S$ is semi-local and $\mathfrak mS$ is contained in the
Jacobson radical of $S$,
\item $M$ is a finite free $S$-module, and
\item $N$ generates $M$ as an $S$-module.
\end{enumerate}
Then $N$ contains an $S$-basis of $M$.
\end{lemma}

\begin{proof}
Assume $M$ is free of rank $n$. Let $I \subset S$ be the Jacobson radical.
By Nakayama's Lemma \ref{lemma-NAK} a sequence of elements
$m_1, \ldots, m_n$ is a basis for $M$ if and only if
$\overline{m}_i \in M/IM$ generate $M/IM$. Hence we may replace
$M$ by $M/IM$, $N$ by $N/(N \cap IM)$, $R$ by $R/\mathfrak m$,
and $S$ by $S/IS$. In this case we see that $S$ is a finite product
of fields $S = k_1 \times \ldots \times k_r$ and
$M = k_1^{\oplus n} \times \ldots \times k_r^{\oplus n}$.
The fact that $N \subset M$ generates $M$ as an $S$-module
means that there exist $x_j \in N$ such that a linear combination
$\sum a_j x_j$ with $a_j \in S$ has a nonzero component in each
factor $k_i^{\oplus n}$.
Because $R = k$ is an infinite field, this means that also
some linear combination $y = \sum c_j x_j$ with $c_j \in k$ has a
nonzero component in each factor. Hence $y \in N$ generates a
free direct summand $Sy \subset M$. By induction on $n$ the result
holds for $M/Sy$ and the submodule $\overline{N} = N/(N \cap Sy)$.
In other words there exist $\overline{y}_2, \ldots, \overline{y}_n$
in $\overline{N}$ which (freely) generate $M/Sy$. Then
$y, y_2, \ldots, y_n$ (freely) generate $M$ and we win.
\end{proof}

\end{document}
